\scanpage{32}
\mspage{\handwritten{11}}
% archive sheet number: 31
\margin{\handwritten{$\unsure{=|--}$}}
\diagram{\begin{tikzcd}[ampersand replacement=\&,column sep=1.5em]
\& \deleted{\unsure{mx|x}}\,P/nP
\handwritten{\arrow[dl,"{x\handwritten{\mapsto}mx}"']}\handwritten{\arrow[dr,"{\handwritten{x\leftrightarrow mx}}"]} \& \\
{}_nP\handwritten{\arrow[dr,"{\deleted{\unsure{x}}\quad x\handwritten{\mapsto x}}"']} \& \&
{}_{\deleted{\unsure{m}}n\handwritten{'}}P\handwritten{\arrow[dl,"{\deleted{\unsure{x}}\ \handwritten{x\mapsto x}}"]}\\
\& {}_{n\handwritten{'}}P \&
\end{tikzcd}}
\deleted{\unsure{ce}} qui \deleted{montre que le diagramme envisagé} est \inserted{\handwritten{\unsure{evidemment|évidemment}}} commutatif.
\margin{\handwritten{$\big|$}}

\deleted{\unsure{2.2.} Soient maintenant}\uline{Corollaire} 2.1\deleted{\unsure{1|x}}.12. Avec les notations de 2.1.11, on a commutativité dans le diagramme suivant
\diagram{\begin{tikzcd}[ampersand replacement=\&,column sep=large,row sep=large]
\& {}_{n\handwritten{'}}P\otimes{}_{n\handwritten{'}}Q
\handwritten{\arrow[r,"{\substack{\handwritten{m\otimes m}\\\deleted{\unsure{p_m\otimes q_m}}}}"]}
\handwritten{\arrow[d,"{\alpha^{n'}_{P,Q}}"']}
\& {}_nP\otimes{}_nQ\handwritten{\arrow[d,"{\alpha^{n}_{P,Q}}"]}\\
\text{(2.1.12.1)} \& {}_{n\handwritten{'}}\underline{\mathrm{Tor}}_1(P,Q)\handwritten{\arrow[r,"{\substack{\handwritten{m}\\\deleted{\unsure{\underline{\mathrm{Tor}}_1(p,q)|\underline{\mathrm{Tor}}_1(P,Q)}}}}"]} \& {}_n\underline{\mathrm{Tor}}_1(P,Q)\ ,
\end{tikzcd}}
\margin{\handwritten{$\xrightarrow{m}$}}
\deleted{où $q_m:{}_{n'}Q\unsure{\longrightarrow|-}{}_nQ$ est \unsure{dééfini|défini} comme $p_m$ comme induit par $m.\mathrm{id}_Q$ . En}
\margin{\handwritten{$\Bigg|$}}
où, dans la description des flèches horizontales, $m$ désigne les homomorphismes $\mathord{}_{n\handwritten{'}}E\handwritten{\longrightarrow}{}_nE$ induits par $\unsure{m|\bar m|\dot m}.\mathrm{id}_E$ (pour $E$ égal respectivement à $P,Q$ et à $\underline{\mathrm{Tor}}_1(P,Q)$).

En effet, en vertu de \deleted{(}2.1.11 on aura, pour $x\handwritten{\in}{}_{n\handwritten{'}}P(S)$, $y\handwritten{\in}{}_{n\handwritten{'}}Q(S)$: $\handwritten{\alpha}^n(mx\handwritten{\otimes}my)=\handwritten{\alpha}^{n\handwritten{'}}(x\handwritten{\otimes}my)$ , or le deuxième membre est égal à $m\handwritten{\alpha}^{n\handwritten{'}}(x\handwritten{\otimes}y)$ p\deleted{\unsure{p}}uisque $x\otimes my=\deleted{\unsure{mx\otimes}}\ m(x\otimes y)$, d'où la commutativité de (2.1.12.1).

2.1.13. Soit $\unsure{l|\ell}$ un nombre premier. Pour tout Groupe commutatif $E$, désignons par $T_l(E)$ le système projectif \deleted{formé des}
\[
T_l(E)=({}_{l^{\handwritten{\nu}}}E)_{\handwritten{\nu\geq}0}\ ,\tag*{(2.1.13.1)}
\]
qui dépend \unsure{évidemment|évidement} fonctoriellement de $E$. Alors les compatibilités 2.1.12, pour $n,n\handwritten{'}$ de la forme $l^{\handwritten{\nu}},l^{\handwritten{\nu'}}$ avec $\handwritten{\nu'\geq\nu}$, nous montrent que les
\margin{\handwritten{$\alpha^{l^\nu}_{P,Q}$}}
\deleted{$\unsure{\alpha|h}$}\inserted{\handwritten{$\alpha^{l^\nu}_{P,Q}$}} pour $\handwritten{\nu}$ variable définissent un homomorphisme de systèmes projectifs\handwritten{,} \inserted{\handwritten{fonctoriel en $P,Q$:}}
\[
\handwritten{\alpha}^{(l)}_{P,Q}\handwritten{:}\ T_l(P)\otimes T_l(Q)\handwritten{\longrightarrow}T_l(\underline{\mathrm{Tor}}_1(P,Q))\ .\tag*{(2.1.13.2)}
\]
La connaissance de ces homomorphismes, pour tout nombre premier $l$, équivaut évidemment à la connaissance des $\handwritten{\alpha}_{P,Q}^n$ pour tout entier $n$. L\handwritten{es}
